% FRAGMENT I01 (integrator): new Online Appendix subsection C.9.
% INSERT BEFORE TEXT: online_appendix.tex "\section{D. Institutional pilot and empirical design}"
% SOURCE: the validated numerical contract, ccc-revision-numerics paper/online_appendix.md,
%   section C.9 (commit 9faf2d7), transcribed into the manuscript's LaTeX conventions.
% NOTATION: the contract's Delta R_T is written \Delta\mathcal R_T, the symbol of (A.41) in the
%   paper for the same object. Displays are unnumbered, as in the contract.
% NUM PLACEHOLDERS (reconciliation names): Pi_acc, Pi_hid, Delta_R_T, break_even_trade_scale,
%   Omega_U, Omega_0, sech_entry_strong, student_t3_entry_strong.
% LABELS ADDED: oa-c-revision (the contract's anchor name).
\subsection{C.9. Gross trade profits, stake bounds, and additional noise
laws}\label{oa-c-revision}

This exercise uses the benchmark with zero low cost. The first
calculation uses \(r_1=3\) and full orders in both price regimes. Write
\(a_H(x)=f(x-1)\), \(a_L(x)=f(x+1)\), and
\(\mu(x)=a_H(x)/(a_H(x)+a_L(x))\). The expected gross trade profit is

\[
\Pi=\Delta_T\int_{\mathbb R}e(x)\frac{a_H(x)a_L(x)}{a_H(x)+a_L(x)}\,dx.
\]

Price access sets \(e(x)=\rho+(1-\rho)\mathbf1\{B_r(\mu(x))\ge c_H\}\).
Hidden prices set \(e(x)=\rho\). Gross profit excludes the trade cost
\(k\). The proceeds gain is the second line of (A.41) in the paper,

\[
\Delta\mathcal R_T=\frac{1-\rho}{2}\left[(t_H-t_0)\alpha_H+(t_L-t_0)\alpha_L\right].
\]

The results are \(\Pi_{\mathrm{acc}}\approx\NUM{Pi_acc}\),
\(\Pi_{\mathrm{hid}}\approx\NUM{Pi_hid}\), and
\(\Delta\mathcal R_T\approx\NUM{Delta_R_T}\). The break-even trade scale
is \(\Delta\mathcal R_T/\Pi_{\mathrm{acc}}\approx\NUM{break_even_trade_scale}\)
order units. These are numerical diagnostics at the stated profile.

The stake calculations use the sufficient bounds of
\hyperref[oa-a-stake]{Section A.12} at the benchmark, where
\(\kappa(1/2)=m\):

\[
\Omega_U=\frac{2b[(1-1/b)\rho m\Delta_T(r_1)-k]}{(1-\rho)(t_H(r_1)-t_0(r_1))},
\qquad
\Omega_0=\frac{2b[k-\Delta_T(r_0)]}{(1-\rho)(t_H(r_0)-t_0(r_0))}.
\]

Their numerical values are \(\Omega_U\approx\NUM{Omega_U}\) and
\(\Omega_0\approx\NUM{Omega_0}\) order units. Use the unrounded formulas
for a strict stake comparison. They do not locate an exact equilibrium
transition.

The additional noise calculations keep full orders fixed at \(r_1\).
For each density, the code solves \(\mu(x)\ge\tau\) and integrates the
two state densities over the complete entry set. Hyperbolic secant noise
gives a half-line. Student noise gives a bounded interval with two
endpoints. The entry probabilities are approximately
\NUM{sech_entry_strong} and \NUM{student_t3_entry_strong}, respectively.
These calculations do not search for equilibria under the additional
laws.

\textbf{Outputs.} \texttt{numerics/revision\_quantities.csv},
\texttt{numerics/revision\_diagnostics.json}, and
\texttt{numerics/manifests/c9\_revision.json}. Each scalar records its
inputs, profile, independent value, discrepancy, and status. Independent
fifty-digit formulas and density integrals validate profits and stakes.
Independent density roots and quadrature validate the noise values. Both
Laplace regimes also pass the existing continuation validator. C.0
controls apply. Any failure leaves the result open and stops the
producer. Feed: Sections 5.1 and 6.1 of the paper and A.12.
